\content{
\begin{example} The baes radius and height of a right circular cone are measured
as 10 cm and 25 cm, respectively, with a possible error in measurement of as
much as 0.01 cm in each. Use a tangent plane to estimate the maximum error in the
calculated volume of the cone. 
\end{example}
}{% presentation
\vspace{4in}
}{% nopresentation
\vspace{2in}
}{% comments
}

\subsection{The Chain Rule}

\content{
\begin{theorem} (The Chain Rule Version 1) Let $f(x,y)$ be a differentiable
function, and $x=g(t)$ and $y=h(t)$ are both differentiable functions. Then 
$$ \frac{df}{dt} = \frac{\partial f}{\partial x}\frac{dx}{dt} + \frac{\partial f}{\partial
y}\frac{dy}{dt}. $$
\end{theorem}
}{% presentation
}{% nopresentation
}{% comments
}

\content{
\begin{example} If $z=x^2y+3xy^4$ where $x=\cos t$ and $y=\sin 3t$, find
$\frac{dz}{dt}$ at $t=0$.
\end{example}
}{% presentation
\vspace{3in}
}{% nopresentation
\vspace{2in}
}{% comments
}

\content{
\begin{theorem} (The Chain Rule Version 2) Let $f(x,y)$ be a differentiable
function where $x=g(s,t)$ and $y=h(s,t)$ are both differentiable function of two
variables, then 
$$ \frac{\partial f}{\partial s} = \frac{\partial f}{\partial x}\frac{\partial x}{\partial s} + 
\frac{\partial f}{\partial y} \frac{\partial y}{\partial s}\quad 
\frac{\partial f}{\partial t} = \frac{\partial f}{\partial x}\frac{\partial x}{\partial t} + 
\frac{\partial f}{\partial y} \frac{\partial y}{\partial t}. $$
\end{theorem}
}{% presentation
}{% nopresentation
\newpage
}{% comments
}

\content{
\begin{example} If $z=e^x\sin y$ and $x=st^2$ and $y=s^2t$, find both 
$ \partial z/\partial s$ and $\partial z/\partial t$.
\end{example}
}{% presentation
\vspace{4in}
}{% nopresentation
\vspace{2in}
}{% comments
}

\content{
\begin{example} If $z = f(x,y)$ has continuous second-order partial derivatives, and $x=s^2+t^2$
and $y=2st$, find expressions for $\partial z/\partial s$ and $\partial^2z/\partial^2 s$.
\end{example}
}{% presentation
\vspace{4in}
}{% nopresentation
\vspace{2in}
}{% comments
}


\content{
\begin{example}
Let $f(x,y) = x\sin(y)$ and suppose that $x = 2st$ and $y=s^3$. Find $\partial
f/
\partial s$ and $\partial f/\partial t$.
\end{example}
}{% presentation
\vspace{3in}
}{% nopresentation
\vspace{2in}
}{% comments
}

\subsubsection{Implicit Differentiation}

\content{
We can use the chain rule to obtain an application to implicit differentiation
from Calculus I:
}{% presentation
\vspace{3in}
}{% nopresentation
\vspace{2in}
}{% comments
Write out $F(x,y) = 0$, and assume that $y=y(x)$, and use the chain rule to
obtain 
$$ \frac{\partial F}{\partial x}\frac{dx}{dx} + \frac{\partial F}{\partial y}\frac{dy}{dx} = 0.
$$
Then since $\frac{dx}{dx}=1$, we can solve for $\frac{dy}{dx}$.
}

\content{
\begin{example} Find $y'$ if $x^3+y^3 = 6xy$.
\end{example}
}{% presentation
\vspace{3in}
}{% nopresentation
\vspace{2in}
}{% comments
}
